% Part: sets-functions-relations
% Chapter: sets
% Section: subsets

\documentclass[../../../include/open-logic-section]{subfiles}

\begin{document}

\olfileid[zh]{sfr}{set}{sub}
\olsection{子集与幂集}

\begin{explain}
我们经常会想要比较集合。而一种显然的比较方式便是：\emph{一个集合中的每个对象也在另一个集合中}。这种情况足够重要，值得我们引入一些新记号。
\end{explain}

\begin{defn}[子集]
若集合 $A$ 的每个!!{element}也是~$B$ 的!!a{element}，则称
$A$ 是~$B$ 的\emph{子集}，记作 $A \subseteq B$。若
$A$ 不是~$B$ 的子集，则记作 $A \not\subseteq B$。
若 $A \subseteq B$ 但 $A \neq B$，则记作 $A \subsetneq B$，并称
$A$ 是 $B$ 的\emph{真子集}。
\end{defn}

\begin{ex}
每个集合都是它自身的子集，且 $\emptyset$ 是每个
集合的子集。自然偶数集合是自然数集的子集。又比如，$\{ a, b \} \subseteq \{ a, b, c \}$。但 $\{ a, b, e
\}$ 不是 $\{ a, b, c \}$ 的子集。
\end{ex}

\begin{ex}
数 $2$ 是整数集的!!{element}，而
偶数集是整数集的子集。然而，一个集合
可能恰好\emph{既}是某个其他集合的!!a{element}，又是它的子集，
例如，$\{0\} \in \{0, \{0\}\}$，且 $\{0\} \subseteq \{0,
\{0\}\}$。
\end{ex}

外延性给出了集合的同一性判据：$A = B$ 当且仅当
$A$ 的每个!!{element}也是~$B$ 的!!a{element}，反之
亦然。「子集」的定义恰恰把 $A \subseteq B$ 界定为
这一判据的前半部分：$A$ 的每个!!{element}也是
~$B$ 的!!a{element}。当然，若我们对调
$A$ 与 $B$，这一定义依然适用：即，$B \subseteq A$ 当且仅当 $B$ 的每个!!{element}也是~$A$ 的!!a{element}。而这，反过来，恰恰就是外延性的
「反之亦然」部分。换言之，外延性
蕴含集合相等当且仅当它们互为子集。

\begin{prop}
$A = B$ 当且仅当既 $A \subseteq B$，又 $B \subseteq A$。
\end{prop}

现在也是再引入一些有用记号的良机。在界定 $A$ 何时是~$B$ 的子集时，我们曾说
「$A$ 的每个!!{element}是 \dots，」并用
「$B$ 的!!a{element}」来填充「$\dots$」。但这是如此常见的一种
\emph{表达形式}，为此引入一些形式化记号将很有帮助。

\begin{defn}\ollabel{forallxina}
$(\forall x \in A)\phi$ 缩写为 $\forall x(x \in A \lif
\phi)$。类似地，$(\exists x \in A)\phi$ 缩写为 $\exists x(x
\in A \land \phi)$。
\end{defn}

借助这一记号，我们可以说 $A \subseteq B$ 当且仅当 $(\forall
x \in A)x \in B$。

接下来我们考虑一类特殊的集合：给定集合的所有
子集的集合。

\begin{defn}[幂集]
由集合~$A$ 的所有子集构成的集合称为
~$A$ 的\emph{幂集}，记作 $\Pow{A}$。
  \[
    \Pow{A} = \Setabs{B}{B \subseteq A}
  \]
\end{defn}

\begin{ex}
$\{ a, b, c \}$ 的所有可能的子集有哪些？它们是：
$\emptyset$，$\{a \}$，$\{b\}$，$\{c\}$，$\{a, b\}$，$\{a, c\}$，$\{b,
c\}$，$\{a, b, c\}$。所有这些子集的集合即
$\Pow{\{a,b,c\}}$：
\[
\Pow{\{ a, b, c \}} = \{\emptyset, \{a \}, \{b\}, \{c\}, \{a, b\},
\{b, c\}, \{a, c\}, \{a, b, c\}\}
\]
\end{ex}

\begin{prob}
列出 $\{a, b, c, d\}$ 的所有子集。
\end{prob}

\begin{prob}
证明：若 $A$ 有 $n$ 个!!{element}，则 $\Pow{A}$ 有 $2^n$
个!!{element}。
\end{prob}

\end{document}
